\documentclass{article}

\textwidth=6.5in
\textheight=8.5in
\voffset=-54pt
\oddsidemargin=0in
\evensidemargin=0in
\setlength{\parskip}{8pt}
\setlength{\parindent}{0pt}

\usepackage{workbook}

\usepackage[sols]{handout}

\begin{document}

\begin{center}
\large\textsc{Problem Set 22 - Definition of the Definite Integral}
  
      \pspace{12pt}
  \begin{problemonly}\large Name: \rule{5cm}{0.15mm}\\ \end{problemonly}
\end{center}


\begin{grading}
\begin{enumerate}
    \item 5 points
    \item 10 points
    \item 20 points
    \item 15 points
    \item 6 points
    \item 4 points
    \item 10 points (EC)
    \item 15 points (EC)
\end{enumerate}
Total: 60 points. (85 points possible with extra credit.)
\end{grading}

\begin{enumerate}

\item 

\begin{grading}
    5 points.
\end{grading}

 \note{Relating Riemann sums and definite
    integrals}
Let $f(x) = \displaystyle \frac{1}{x}$.  Suppose we want to
  approximate $\displaystyle \int_1^3 f(x)\,dx$ using left-hand sums $L_n$
  and right-hand sums $R_n$.

    \begin{problem}[\vfill]
      Put the following nine expressions in ascending order.  (Put ``$<$'' or
      ``$=$'' signs between them.)
      \[ 
      \int_1^3 f(x) \, dx, \quad L_5, \quad R_5, \quad L_{20}, \quad
      R_{20}, \quad L_{100}, \quad R_{100}, \quad \lim_{n \to \infty} L_n,
      \quad \lim_{n \to \infty} R_n\]
    \end{problem}

    \begin{solution}
      (This problem is very similar to
          {{{{\#4}} on {the worksheet ``The Definite Integral''}}} and
          {{{{\#5}}(a)}} on the worksheet.)

      Since $\frac{1}{x}$ is decreasing on $[1, 3]$, left-hand sums
      will be overestimates for $\displaystyle \int_1^3 f(x)\,dx$, and
      right-hand sums will be underestimates.  Both $\displaystyle \lim_{n
        \to \infty} L_n$ and $\displaystyle \lim_{n \to \infty} R_n$ are
      exactly equal to $\displaystyle \int_0^2 f(x)\,dx$, by the
      definition of the definite integral. 

      Here are pictures of $L_5$ and $L_{20}$:
      \begin{center}
        \begin{tikzpicture}[declare function={f(\x)=1/\x;}]
          \begin{axis}[ytick=\empty, xtick={1, 3}, cycle list=black,
            clip=false, hide y axis] 
            \addplot+[ybar interval, plusarea] coordinates {
              (1, {f(1)}) (7/5, {f(7/5)}) (9/5, {f(9/5)}) (11/5,
              {f(11/5)}) (13/5, {f(13/5)}) (3, {f(13/5)}) }; 
            \addplot+[domain=0.9:3.1] {f(x)};
            \draw (axis cs: 0.5, 0.5) node {$L_5$};
          \end{axis}
        \end{tikzpicture}

        \begin{tikzpicture}[declare function={f(\x)=1/\x;}]
          \begin{axis}[ytick=\empty, xtick={1, 3}, cycle list=black,
            clip=false, hide y axis] 
            \addplot+[ybar interval, plusarea] coordinates {
              (1, {f(1)}) (11/10, {f(11/10)}) (6/5, {f(6/5)}) (13/10,
              {f(13/10)}) (7/5, {f(7/5)}) (3/2, {f(3/2)}) (8/5, {f(8/5)})
              (17/10, {f(17/10)}) (9/5, {f(9/5)}) (19/10, {f(19/10)}) (2,
              {f(2)}) (21/10, {f(21/10)}) (11/5, {f(11/5)}) (23/10,
              {f(23/10)}) (12/5, {f(12/5)}) (5/2, {f(5/2)}) (13/5,
              {f(13/5)}) (27/10, {f(27/10)}) (14/5, {f(14/5)}) (29/10,
              {f(29/10)}) (3, {f(29/10)})
              };
            \addplot+[domain=0.9:3.1] {f(x)};
            \draw (axis cs: 0.5, 0.5) node {$L_{20}$};
          \end{axis}
        \end{tikzpicture}
      \end{center}
      As we can see, $L_{20}$ is much closer to $\displaystyle \int_1^3
      f(x)\,dx$ than $L_5$ is.  By the same idea, $L_{100}$ will be even
      closer to $\displaystyle \int_1^3 f(x)\,dx$ than $L_{20}$ is.
      Similarly, among the right-hand sums, $R_{100}$ is closest to
      $\displaystyle \int_1^3 f(x)\,dx$, followed by $R_{20}$ and then
      finally by $R_5$.  That gives us the following ordering:
      \[ \boxed{R_5 < R_{20} < R_{100} < \int_1^3 f(x)\,dx =
        \lim_{n \to \infty} L_n = \lim_{n \to \infty} R_n < L_{100} <
        L_{20} < L_5}.\]
    \end{solution}



\item 

\begin{grading}
10 points. 3 for each part, and 1 for comparing their answers.
\end{grading}

We've defined the definite integral as a limit of Riemann sums and seen
  that it is more efficient to write Riemann sums in $\sum$ notation.  In
  this problem, we just want to make sure you have some practice
  ``unpacking'' $\sum$ notation.  (You can read more about $\sum$ notation
  in Gottlieb \S 18.4; ignore the references to geometric series.)

  For each of the following sums write out the first three terms ``$+ \cdots + $'' the last
  term.  For example, $\displaystyle \sum_{k=1}^{100} k = 1 + 2 + 3 +
  \cdots + 100$.

  \emph{(If you plan on taking Math 1b, being comfortable with this
    notation will be very important!)}

\begin{enumerate}
    % \item $\displaystyle \sum_{i=2}^5 i^2$

    %   \begin{solution}
    %     $\displaystyle \sum_{i=2}^5 i^2 = 2^2 + 3^2 + 4^2 + 5^2 =
    %     4+9+16+25 = 54$ 
    %   \end{solution}

    \item
\begin{problem}[\vfill]
        $\displaystyle \sum_{k=1}^{200} f(x_k) \Delta x$
      \end{problem}

      \begin{solution}
        $\displaystyle \sum_{k=1}^{200} f(x_k) \Delta x = f(x_1)\Delta x +
        f(x_2)\Delta x + f(x_3)\Delta x + \cdots + f(x_{199})\Delta x +
        f(x_{200})\Delta x$. 
      \end{solution}

    \item
\begin{problem}[\vfill]
        $\displaystyle \sum_{k=1}^{200} f(x_{k-1}) \Delta x$        
      \end{problem}

      \begin{solution}
        $\displaystyle \sum_{k=1}^{200} f(x_{k-1}) \Delta x = f(x_0)
        \Delta x + f(x_1) \Delta x + f(x_2) \Delta x + \cdots + f(x_{198})
        \Delta x + f(x_{199}) \Delta x$.        
      \end{solution}

    \item
\begin{problem}[\vfill]
        $\displaystyle \sum_{k=0}^{199} f(x_k) \Delta x$        
      \end{problem}

      \begin{solution}
        $\displaystyle \sum_{k=0}^{199} f(x_k) \Delta x = f(x_0) \Delta x
        + f(x_1) \Delta x + f(x_2) \Delta x + \cdots + f(x_{198}) \Delta x
        + f(x_{199}) \Delta x$.
      \end{solution}

%     \item
% \begin{problem}
%         $\displaystyle \sum_{k=1}^n \left( \frac{k}{n} \right)^2$
%       \end{problem}

%       \begin{solution}
%         $\displaystyle \sum_{k=1}^n \left( \frac{k}{n} \right)^2 =
%         \left(\dfrac{1}{n}\right)^2 + \left(\dfrac{2}{n}\right)^2 + 
% 	\left(\dfrac{3}{n}\right)^2 + \cdots
%         +\left(\dfrac{n-1}{n}\right)^2 + \left(\dfrac{n}{n}\right)^2 $. 
%       \end{solution}

%     \item
% \begin{problem}
%         $\displaystyle \sum_{k=1}^n \left( \frac{k-1}{n} \right)^2$
%       \end{problem}

%       \begin{solution}
%         $\displaystyle \sum_{k=1}^n \left( \frac{k-1}{n} \right)^2 =
%         \left(\dfrac{0}{n}\right)^2 + \left(\dfrac{1}{n}\right)^2 +
%         \left(\dfrac{2}{n}\right)^2 + \cdots
%         +\left(\dfrac{n-2}{n}\right)^2 + \left(\dfrac{n-1}{n}\right)^2$.
%       \end{solution}

%     \item
% \begin{problem}
%         $\displaystyle \sum_{k=0}^{n-1} \left( \frac{k}{n} \right)^2$
%       \end{problem}

%       \begin{solution}
%         $\displaystyle \sum_{k=0}^{n-1} \left( \frac{k}{n} \right)^2 =
%         \left(\dfrac{0}{n}\right)^2 + \left(\dfrac{1}{n}\right)^2 +
%         \left(\dfrac{2}{n}\right)^2 + \dots +\left(\dfrac{n-2}{n}\right)^2
%         + \left(\dfrac{n-1}{n}\right)^2$.
%       \end{solution}

\end{enumerate}

    \begin{problem}[\vfill\newpage]
      Compare your three answers. Write 1-2 sentences about
      similarities you see.

      % How about {{{{\#3}}(e)}},
      % {{{{\#3}}(f)}}, and
      % {{{{\#3}}(g)}}?
    \end{problem}

    \begin{solution}
      {{{{\#2}}(b)}} and
      {{{{\#2}}(c)}} are the same.
      {{{{\#2}}(a)}} is similar: the pattern of
      terms is the same, but it starts and ends in slightly different
      places. 

      % Similarly, {{{{\#3}}(f)}} and
      % {{{{\#3}}(g)}} are the same.
    \end{solution}



\item 

\begin{grading}
20 points, as follows: 2/3/2/3/2/3/5
\end{grading}

{%
    A health clinic has set up a flu shot program for its patients.  The
    clinic is open from 8 am to 4 pm (16:00) and gives flu shots on a
    first-come, first-serve basis. The clinic has the capacity to serve 30
    patients per hour. The function $r(t)$, whose graph is given below,
    gives the rate at which people are arriving at the clinic for shots
    one day.

    \providecommand{\graphcommand}{}
    \renewcommand{\graphcommand}[1]{
      \begin{tikzpicture}[declare function={f(\x) = 69 - 3*sqrt(4369 -
          720*\x + 30*\x ^2);}]
        \begin{axis}[cycle list=black, xtick={8, 9, ..., 16}, ytick={15, 30,
            45}, grid=major, clip=false, enlargelimits=false, width=2.5in,
          height=1.5in, axis line style={shorten >=-8pt}, xlabel
          style={xshift=8pt}, ylabel style={yshift=8pt}, 
          xlabel={$t$ (time)},
          ylabel={$r(t)$ (people/hr)}]
          #1
          \addplot+[domain=8:16] {f(x)};
        \end{axis}
      \end{tikzpicture}
    }
    \graphcommand{}

    Explain your answers to the questions below by relating them to points,
    lengths, or areas on the graph.

    \begin{enumerate}
    
    \item
      \begin{problem}[\vfill]
        Approximate the number of people who came to the clinic for a flu
        shot this day.
      \end{problem}

      \begin{solution} % Jason Bland
        This is the entire area under $r(t)$ from $t = 8$ to $t = 16$.
        This is about 15 grid squares in all, or about \fbox{225 people}.
      \end{solution}

    \item 
      \begin{problem}[\vfill]
        At what time does a line start forming? Explain why there is no
        line at 9 am even though people have been coming to the clinic
        since 8 am.
      \end{problem}

      \begin{solution} % Jason Bland
        This happens once people begin to arrive at the faster rate than
        the clinic can serve (30 people per hour), which is at \fbox{10
          am}.
      \end{solution}

    \item
\begin{problem}[\vfill]
        At approximately what time is the length of the line increasing most
        rapidly?
      \end{problem}

      \begin{solution} % Jason Bland
        This is when people are arriving at the fastest rate, which is
        at \fbox{noon}.
      \end{solution}

    \item
\begin{problem}[\nstretch{2}]
        At approximately what time is the line the longest? The answer to
        this is different from the answer to part (c); explain
        why. (Hint: the answer is NOT $12$ and NOT $16$.)
      \end{problem}

      \begin{solution} % Jason Bland
        Between 10 am and 2 pm, the length of the line is increasing; after
        2 pm, the length of the line is decreasing (the rate at which people
        arrive is less than the rate at which they are served).  Therefore,
        the line is longest \fbox{at 2 pm ($t = 14$)}.

        This differs from {{{{\#5}}(b)}} because, between noon and 2
        pm, people are arriving at a rate which is slower than the rate at
        noon but still higher than 30 per hour. This means the line will
        continue to lengthen, just not as quickly as at noon.
      \end{solution}

    \item
\begin{problem}[\vfill]
        When the line is longest, approximately how many people are in line?
      \end{problem}

      \begin{solution} 
        We're asked to find the length of the line at 2 pm ($t = 14$).
        Since the line started forming at $t = 10$, we can think of this
        as (number of people who arrived between $t = 10$ and $t = 14$)
        $-$ (number of people who were served between $t = 10$ and $t =
        14$).  The number of people who arrived between $t = 10$ and $t =
        14$ can be pictured as the area under $r(t)$ from $t = 10$ to $t =
        14$ (the area marked in blue below); this area is about 165 (about
        11 grid squares).  The number of people who were served between $t
        = 10$ and $t = 14$ was (30 people / hour)(4 hours) = 120, so there
        are approximately $165 - 120 = \boxed{45}$ people in line.

        Here's an alternate way to picture it: the number of people who
        were served between $t = 10$ and $t = 14$ can be pictured as the
        area under $y = 30$ from $t = 10$ to $t = 14$ (marked in red
        below).  Then, the length of the line is the area of the region
        which is just blue (not red), which is about 45 people (3 grid
        squares): 
        \begin{center}
          \graphcommand{
            \addplot+[domain=10:14, fill=cyan, fill opacity=0.5, draw=none] {f(x)} \closedcycle;
            \addplot+[domain=10:14, pattern=flexible hatch, pattern
            color=red, draw=none] {30} \closedcycle;
            \addplot+[domain=10:14, red] {30};
          }
        \end{center}
      \end{solution}

    \item
      \begin{problem}[\vfill]
        Approximate the longest amount of time a person could wait for a
        shot. (Assume that doors close to new arrivals at 4 pm ($t=16$) but 
        everyone who has arrived by 4 pm is served.)
      \end{problem}

      \begin{solution} 
        The amount of time a person waits depends on how many people are
        in line before him.  We already determined that there are about 45
        people in line when the line is longest.  A person arriving at
        this time has to wait for those 45 people to be served, which
        takes $\frac{45 \text{ people}}{30 \text{ people / hr}} =
        \boxed{1.5 \text{ hours}}$.
      \end{solution}

    \item
      \begin{problem}[\vfill\newpage]
        Approximately how long is the line at 3 pm?
      \end{problem}

      \begin{solution} % Jason Bland
        We already figured out that the line has about 45 people in it at
        2 pm, so let's figure out what happens between 2 and 3 pm.  In
        this time, 30 people are served, but roughly 23 people arrive
        (this is the area under $r(t)$ from $t = 14$ to $t = 15$, which is
        about 1.5 grid squares).  So, at 3 pm, the line has about $45 - 30
        + 23 = \boxed{38 \text{ people}}$.
      \end{solution}

    \end{enumerate}
  }

\item 

\begin{grading}
15 points, as follows: 1/2/2/2/4/4
\end{grading}

For the scenario described in {{\#3}}, express the
  following quantities exactly; your answers may be in terms of definite
  integrals.  (For instance, your answer might look like $50 +
  \displaystyle \int_8^{10} r(t)\,dt$ or $\displaystyle \int_8^9 [r(t) -
  2]\,dt$.)

  \begin{enumerate}
    \item
    \begin{problem}[\vfill]
      The total number of people who came to the clinic for flu shots this day 
    \end{problem}

    \begin{solution}
      As we saw in {{{{\#3}}(a)}}, this is the entire area $r(t)$
      from $t = 8$ to $t = 16$, or $\boxed{\int_8^{16} r(t) \, dt}$.
    \end{solution}
    
  \item
\begin{problem}[\vfill]
      The length of the line at 11 am
    \end{problem}

    \begin{solution} 
      Using the same reasoning as in {{{{\#3}}(e)}}, the length of
      the line at 11 am is (number of people who arrive between $t = 10$
      and $t = 11$) $-$ (number of people served between $t = 10$ and $t =
      11$), or $\boxed{\int_{10}^{11} r(t)\,dt - 30}$.
    \end{solution}

  \item
\begin{problem}[\vfill]
      The length of the line at its longest
    \end{problem}

    \begin{solution} 
      We saw in {{{{\#3}}(d)}} that the line is longest at 2 pm,
      and the length of the line at 2 pm is $\boxed{\int_{10}^{14}
        r(t)\,dt - 120}$.  (Here, $\displaystyle \int_{10}^{14} r(t)\,dt$
      is the number of people who arrived between 10 am and 2 pm, and
      $120$ is the number of people served between 10 am and 2 pm.)
    \end{solution}

  \item
    \begin{problem}[\vfill]
      The length of the line at 3 pm
    \end{problem}

    \begin{solution} % Jason Bland
      This is essentially the same as {{{{\#4}}(b)}} and
      {{{{\#4}}(c)}}; the number of people who arrive between 10
      am and 3 pm is $\displaystyle \int_{10}^{15} r(t)\,dt$, while the
      number of people served in that time is (30 people / hour)(5 hours)
      = 150.  So, the length of the line at 3 pm is $\boxed{\int_{10}^{15}
        r(t)\,dt - 150}$.
    \end{solution}

  \item
    \begin{problem}[\vfill]
      The amount of time a person arriving at noon has to wait in line
    \end{problem}

    \begin{solution} % Jason Bland
      The length of the line at noon is $\displaystyle \int_{10}^{12}
      r(t)\,dt - 60$.  A person arriving at noon has to wait for all of
      these people to be served, which happens at a rate of 30 people per
      hour.  So, the waiting time is $\boxed{\frac{\displaystyle
          \int_{10}^{12} r(t)\,dt - 60 \text{ people}}{30 \text{
            people/hour}}}$ (the units simplify to hours).
    \end{solution}

  \item
    \begin{problem}[\vfill\newpage]
      The amount of time by which the clinic must extend its hours in
      order to serve everyone who arrives by 4pm
    \end{problem}

    \begin{solution} 
      At 4 pm, the number of people in line is $\displaystyle
      \int_{10}^{16} r(t)\,dt - 180$.  The amount of time it would take to
      serve these people is $\boxed{\frac{\displaystyle
      \int_{10}^{16} r(t)\,dt - 180 \text{ people}}{30 \text{
        people/hour}}}$. 
    \end{solution}

  \end{enumerate}

% \item

%  (22.2.5)
% Below are the velocity graphs for a chicken and a goat. Assume that at
%   time $t=0$ the goat and the chicken start out side by side and they both
%   travel along the same straight dirt path. Answer the questions below.
%   When the quantity corresponds to an area, describe this area on the
%   graph and then give it using the appropriate definite integral or sums
%   and differences of definite integrals. The graphs intersect at $t=1.5$
%   and $t=6.5$. The graph of $v_g(t)$ has a maximum at $t=4.5$.

%   \includegraphics{images/go22m198}

%   \begin{enumerate}

%   \item
%     \begin{problem}
%       The lead the chicken has when the goat begins to move
%     \end{problem}

%     \begin{solution}
%       The goat begins to move at $t = 1$, so the chicken's lead is its net
%       change in position from $t = 0$ to $t = 1$, or $\boxed{\int_0^1
%         v_c(t)\, dt}$.
%     \end{solution}

%   \item
%     \begin{problem}
%       The distance between the goat and the chicken at time $t=1.5$.
%     \end{problem}

%     \begin{solution}
%       The distance the chicken has traveled by $t = 1.5$ is $\displaystyle
%       \int_0^{1.5} v_c(t)\,dt$, while the distance the goat has traveled
%       is $\displaystyle \int_0^{1.5} v_g(t)\,dt$.  Therefore, the distance
%       between the two is $\boxed{\int_0^{1.5} v_c(t)\,dt - \int_0^{1.5}
%         v_g(t)\,dt}$.  (We can be sure that the chicken is ahead because,
%       on $[0, 1.5]$, it's always moving faster than the goat.)
%     \end{solution}

%   \item
%     \begin{problem}
%       The distance between the goat and the chicken at time $t=3$
%     \end{problem}

%     \begin{solution}
%       The distance the chicken has traveled by $t = 3$ is $\displaystyle
%       \int_0^{3} v_c(t)\,dt$, while the distance the goat has traveled
%       is $\displaystyle \int_0^{3} v_g(t)\,dt$.  It's not totally clear
%       which one is ahead at $t = 3$, so we should use absolute values to
%       describe the distance between them: $\boxed{\left| \int_0^3
%           v_c(t)\,dt - \int_0^3 v_g(t)\,dt \right|}$. 
%     \end{solution}

%   \item
%     \begin{problem}
%       The time at which the goat is farthest ahead of the chicken
%     \end{problem}

%     \begin{solution}
%       The goat is farthest ahead of the chicken at $\boxed{t=6.5}$.
%       Starting at $t = 6.5$, the chicken moves faster than the goat and
%       therefore starts to catch up to the goat. 
%     \end{solution}

%   \item
%     \begin{problem}
%       Approximate the time(s) at which the chicken and the goat pass one
%       another on the path.
%     \end{problem}

%     \begin{solution}
%       These are the times at which the chicken and goat have traveled
%       the same distance; that is, they are the times $T$ when
%       $\displaystyle \int_0^T v_c(t)\,dt = \int_0^T v_g(t)\,dt$.  It looks
%       like the only such time is \fbox{around $3$}.
%     \end{solution}

%   \item
%     \begin{problem}
%       Approximate the time at which the distance between the goat and the
%       chicken is increasing most rapidly.
%     \end{problem}

%     \begin{solution}
%       This is the time at which the two animals' velocities differ the
%       most.  This appears to be \fbox{around $t = 4.25$}. 
%     \end{solution}

%   \end{enumerate}

\item 

\begin{grading}
6 points. 3 for each part.
\end{grading}

 % 22.4.3
  For each of the two integrals below, choose the correct answer and explain your reasoning.  (\emph{Hint:}
  Think about symmetry of the integrand.\footnote{In a definite integral
    $\displaystyle \int_a^b f(t)\,dt$, the \uline{integrand} is the
    function $f(t)$.})
  \begin{enumerate}

  \item $\displaystyle\int_{-a}^{a} \frac{1}{1+x^2}\,dx =$

    \begin{tabular*}{5.5in}{@{\extracolsep{\fill}}cccc}
        i.  $0$ & ii.  $\displaystyle \int_{0}^{a}\frac{1}{1+x^2}\,dx $
        & iii. \probonly{$\displaystyle 2\int_{0}^{a} \frac{1}{1+x^2}\,dx$}
        \solonly{$\boxed{\displaystyle 2\int_{0}^{a} \frac{1}{1+x^2}\,dx}$}
        & iv.  $\displaystyle -\int_{0}^{a}\frac{1}{1+x^2}\,dx$
    \end{tabular*}

    \begin{solution}
      Since $\frac{1}{1 + x^2}$ is an even function, the signed area under it
      from $x = -a$ to $x = a$ should be twice the signed area from $0$ to $a$. 
    \end{solution}

    \pspace{\vfill}

  \item $\displaystyle\int_{-a}^{a}\frac{x}{1+x^2}\,dx =$

    \begin{tabular*}{5.5in}{@{\extracolsep{\fill}}cccc}
        i.  \probonly{$0$} \solonly{$\boxed{0}$} 
        & ii.  $\displaystyle \int_{0}^{a}\frac{x}{1+x^2}\,dx$
        & iii. $\displaystyle 2\int_{0}^{a} \frac{x}{1+x^2}\,dx$
        & iv.  $\displaystyle -\int_{0}^{a}\frac{x}{1+x^2}\,dx$
    \end{tabular*}
    

    \begin{solution}
      Since $\frac{x}{1 + x^2}$ is an odd function, the signed area under
      it from $x = -a$ to $x = a$ is 0. 
    \end{solution}

    \pspace{\vfill}

  \end{enumerate}

\item 

\begin{grading}
4 points.
\end{grading}

\begin{problem}
    \textbf{Reflect Back}.  Suppose we use right-hand and left-hand sums
    to estimate the areas pictured below.  Match the pictures with the
    statements. Note that there is a one-to-one correspondence, so every picture should be matched with a different statement.
    \begin{center}
      \begin{tabular}{c c c c}
        (a) & (b) & (c) & (d) \\
        \includegraphics[width=\linewidth/4-\tabcolsep*2]{images/area-sums-reflect-back-3}
        & 
        \includegraphics[width=\linewidth/4-\tabcolsep*2]{images/area-sums-reflect-back-1}
        & 
        \includegraphics[width=\linewidth/4-\tabcolsep*2]{images/area-sums-reflect-back-4}
        & 
        \includegraphics[width=\linewidth/4-\tabcolsep*2]{images/area-sums-reflect-back-2}
        \\
      \end{tabular}
    \end{center}

    Choices:
    \begin{enumerate}[label=\Alph*.]
    \item $L_4$ will give an underestimate for this area, $R_4$ an
      overestimate.
    \item $L_4$ will give an overestimate, $R_4$ an underestimate.
    \item It is very hard to tell from the picture whether $L_4$ and $R_4$
      will be overestimates or underestimates.
    \item Both $L_4$ and $R_4$ will give overestimates.
    \end{enumerate}
  \end{problem}

  \begin{solution} % Jason Bland
    (a) goes with (B), (b) goes with (A), (c) goes with (D), and (d) goes with (C).
  \end{solution}

  \pspace{\vfill\newpage}


%   \item 

%   \begin{grading}
% 10 points, as follows: 5/2/3.
% \end{grading}

%   \textbf{Extra Credit:} As in problem \#1, let $f(x) = \displaystyle \frac{1}{x}$, and suppose we want to approximate $\displaystyle \int_1^3 f(x)\,dx$.  In this problem, you will figure out how many approximating rectangles we need to use in order to estimate the integral within a certain margin of error. (This type of problem is called an ``error-bounding" problem, and they occur frequently in Math 1b!)
  
% \begin{enumerate}
%   \item
% \begin{problem}[\vfill]
%       Find $|R_5 - L_5|$.  (\emph{Hint:} You don't need to calculate $L_5$
%       or $R_5$ to do this.  If you're puzzled, try writing out expressions
%       for $R_5$ and $L_5$ and/or looking at $L_5$ and $R_5$ in the Riemann
%       sums applet.  If you're still stuck, take a look at the solution to
%       {{{{\#5}}(b)}} on the worksheet.) 
%     \end{problem}

%     \begin{solution}
%       (This is very similar to
%       {{{{\#5}}(b)}} on the worksheet.)

%       Here are pictures of $L_5$ (in red stripes) and $R_5$ (in blue
%       stripes): 
%       \begin{center}
%         \begin{tikzpicture}[declare function={f(\x)=1/\x;}]
%           \begin{axis}[ytick=\empty, xtick={1, 3}, cycle list=black,
%             clip=false, hide y axis, xtick={1, 1.4, ..., 3.01},
%             xticklabels={$x_0$, $x_1$, $x_2$, $x_3$, $x_4$, $x_5$},
%             ymin=0]            
%             \addplot+[ybar interval, pattern=flexible hatch, pattern
%             color=red] coordinates { (1, {f(1)}) (7/5, {f(7/5)}) (9/5,
%               {f(9/5)}) (11/5, {f(11/5)}) (13/5, {f(13/5)}) (3, {f(13/5)})
%             }; %
%             \addplot+[ybar interval, pattern=flexible hatch, pattern
%             color=blue, hatch direction=-1] coordinates { (1, {f(7/5)})
%               (7/5, {f(9/5)}) (9/5, {f(11/5)}) (11/5, {f(13/5)}) (13/5,
%               {f(3)}) (3, {f(3)}) }; %
%             \addplot+[domain=0.9:3.1] {f(x)}; %
%             \draw (axis cs: 1, 0) node[font=\footnotesize, yshift=-12pt,
%             rotate=90] {$=$}; %
%             \draw (axis cs: 1, 0) node[font=\footnotesize, yshift=-20pt]
%             {$1$}; %
%             \draw (axis cs: 3, 0) node[font=\footnotesize, yshift=-12pt,
%             rotate=90] {$=$}; %
%             \draw (axis cs: 3, 0) node[font=\footnotesize, yshift=-20pt]
%             {$3$}; %
%           \end{axis}
%         \end{tikzpicture}
%       \end{center}
%       Notice that the first 4 rectangles in $R_5$ are the same as the last
%       4 rectangles in $L_5$; that is,
%       \[ R_5 - \text{(last
%         rectangle in $R_5$)} = L_5 - \text{(first rectangle in $L_5$)}.\]
%       Rearranging this equation,
%       \begin{align*}
%         R_5 - L_5 &= \text{(last rectangle in $R_5$)} - \text{(first
%           rectangle in $L_5$)} \\
%         &= f(3) \Delta x - f(1) \Delta x \text{ where $\Delta x =
%           \frac{2}{5}$ is the width of each rectangle} 
%       \end{align*}
%       So, $|R_5 - L_5| = |f(3) - f(1)| \Delta x = \left| \frac{1}{3} - 1
%       \right| \frac{2}{5} = \boxed{\frac{4}{15}}$. 

%       If you prefer, you could write out both $L_5$ and $R_5$:
%       \begin{align*}
%         R_5 &= f(x_1) \Delta x + f(x_2) \Delta x + f(x_3) \Delta x +
%         f(x_4) \Delta x + f(x_5) \Delta x \\
%         L_5 &= f(x_0) \Delta x + f(x_1) \Delta x + f(x_2) \Delta x +
%         f(x_3) \Delta x + f(x_4) \Delta x
%       \end{align*}
%       Therefore, $R_5 - L_5 = f(x_5) \Delta x - f(x_0) \Delta x = f(3)
%       \Delta x - f(1) \Delta x$, as we said before.
%     \end{solution}

%   \item
%     \begin{problem}[\vfill]
%       Find $|R_{100} - L_{100}|$.
%     \end{problem}

%     \begin{solution}
%       By the same reasoning as in
%       {{{{\#7}}(a)}}, $|R_{100} -
%       L_{100}| = |f(3) - f(1)| \Delta x$ where $\Delta x = \frac{2}{100}$
%       (this is the width of each rectangle if we use 100 rectangles to
%       approximate $\displaystyle \int_1^3 f(x)\,dx$).  That is, $|R_{100}
%       - L_{100}| = \left| \frac{1}{3} - 1 \right| \frac{2}{100} =
%       \frac{4}{300} = \boxed{\frac{1}{75}}$. 
%     \end{solution}

%   \item
%     \begin{problem}[\vfill]
%       How large must $n$ be to ensure that $|R_n-L_n| < 0.01$?
%     \end{problem}

%     \begin{solution}
%       By the same reasoning as in
%       {{{{\#7}}(a)}}, $|R_n - L_n|
%       = |f(3) - f(1)| \Delta x$ where $\Delta x = \frac{2}{n}$ (this is
%       the width of each rectangle if we use $n$ rectangles to approximate
%       $\displaystyle \int_1^3 f(x)\,dx$).  That is, $|R_n - L_n| = \left|
%         \frac{1}{3} - 1 \right| \frac{2}{n} = \frac{4}{3n}$.  We want this
%       to be less than $0.01$; that is, we want
%       \begin{align*}
%         \frac{4}{3n} &< 0.01 \\
%         \frac{4}{3 (0.01)} & < n \\
%         \frac{400}{3} & < n 
%       \end{align*}
%       So, we'd like $n$ to be \fbox{at least 134} to be sure that $|R_n -
%       L_n| < 0.01$.
%     \end{solution}

%     \probonly{\emph{Knowing this enables us to pin $\displaystyle \int_1^3
%         f(x)\,dx$ between two numbers that are no more than $0.01$ apart,
%         meaning we can approximate $\displaystyle \int_1^3 f(x)\,dx$ to
%         within $0.01$.}}

%         \pspace{\newpage}
%   \end{enumerate}

  

\item 
\begin{grading}
15 points, as follows: 3/2/2/3/4/1
\end{grading}

  In this problem, you'll look at $\displaystyle \int_0^1 x^3\,dx$, and compute its exact value using Riemann sums. (See \#4 in workbook section 5.3 for a similar problem from class.)
  \begin{enumerate}
  \item
\begin{problem}[\vfill]
      Write out $R_4$, the right-hand sum with 4 rectangles that
      approximates this integral.  Evaluate $R_4$.  (Use the Riemann Sums
      applet at \url{https://www.geogebra.org/m/sgxnjavg} to check
      your answer.)
    \end{problem}

    \begin{solution}
      We're looking at the interval $[0, 1]$, which has width 1.  So, if
      we use 4 rectangles, each rectangle should have width
      $\frac{1}{4}$.  Here's a picture:
      \begin{center}
        \begin{tikzpicture}[declare function={f(\x) = \x^3;}]
          \begin{axis}[cycle list=black, xtick={0, 0.25, ..., 1.01},
            xticklabels={$0$, $\frac{1}{4}$, $\frac{2}{4}$, $\frac{3}{4}$,
              $1$}, ytick=\empty]
            \addplot+[ybar interval, plusarea] coordinates { (0, {f(1/4)})
              (1/4, {f(1/2)}) (1/2, {f(3/4)}) (3/4, {f(1)}) (1, {f(1)}) };
            \addplot+[domain=0:1.05] {f(x)} node[left]{$y = x^3$};
          \end{axis}
        \end{tikzpicture}
      \end{center}
      So, $\displaystyle R_4 = \left( \frac{1}{4} \right)^3 \frac{1}{4} +
      \left( \frac{2}{4} \right)^3 \frac{1}{4} + \left( \frac{3}{4}
      \right)^3 \frac{1}{4} + \left( \frac{4}{4} \right)^3 \frac{1}{4} =
      \boxed{\frac{25}{64}}$.
    \end{solution}

  \item
\begin{problem}[\vfill]
      Write out $R_{50}$, but do not evaluate.  ($R_{50}$ involves adding
      $50$ terms; we don't want you to write them all out!  Instead, write
      out the first three terms $+ \cdots + $ and then the last two
      terms.)
    \end{problem}

    \begin{solution}
      Using the same reasoning as in
      {{{{\#8}}(a)}}, \[R_{50} =
      \boxed{\left( \frac{1}{50} \right)^3 \frac{1}{50} + \left(
          \frac{2}{50} \right)^3 \frac{1}{50} + \left( \frac{3}{50}
        \right)^3 \frac{1}{50} + \cdots + \left( \frac{49}{50} \right)^3
        \frac{1}{50} + \left( \frac{50}{50} \right)^3 \frac{1}{50}}.\]
    \end{solution}

  \item
    \begin{problem}[\nstretch{3}]
      It turns out that there is a formula for the sum $1^3+2^3+3^3+\cdots+n^3$, shown below.
      \begin{equation*}
        1^3 + 2^3 + 3^3 + \cdots + n^3 = \dfrac{n^4 + 2n^3 +
          n^2}{4}
      \end{equation*}
      When $n = 50$, $1^3 + 2^3 + 3^3 + \cdots + 50^3 =
      \dfrac{50^4 + 2 \cdot 50^3 + 50^2}{4} = 1,625,625$.  \\ Use this to
      evaluate $R_{50}$.  (Use the Riemann Sums applet to check your
      answer.)
    \end{problem}

    \begin{solution}
      From {{{{\#8}}(b)}}, 
      \begin{align*}
        R_{50} & = \left( \frac{1}{50} \right)^3 \frac{1}{50} + \left(
          \frac{2}{50} \right)^3 \frac{1}{50} + \left( \frac{3}{50}
        \right)^3 \frac{1}{50} + \cdots + \left( \frac{49}{50} \right)^3
        \frac{1}{50} + \left( \frac{50}{50} \right)^3 \frac{1}{50}\\
        &= \frac{1}{50^4} + \frac{2^3}{50^4} + \frac{3^3}{50^4} + \cdots +
        \frac{49^3}{50^4} + \frac{50^3}{50^4} \\
        & = \frac{1}{50^4}\left(1^3 + 2^3 + 3^3 + \cdots + 50^3  \right) \\
        & = \frac{1}{50^4}(1625625) \text{ by (*)} \\
        & = \boxed{\frac{2601}{10000}}
      \end{align*}
    \end{solution}

  \item
    \begin{problem}[\stretch{2.5}]
      Use the formula on the previous page for $1^3+2^3+3^3+\cdots+n^3$ to find a  formula for $R_n$.

      \checkmark\ Check your formula by plugging in $n = 4$ and $n = 50$.
    \end{problem}

    \begin{solution}
      Using the same reasoning as in
      {{{{\#8}}(a)}}, 
      \begin{align*}
        R_{n} & = \left( \frac{1}{n} \right)^3 \frac{1}{n} + \left( \frac{2}{n}
        \right)^3 \frac{1}{n} + \left( \frac{3}{n} \right)^3
        \frac{1}{n} + \cdots + \left( \frac{n-1}{n} \right)^3
        \frac{1}{n} + \left( \frac{n}{n} \right)^3 \frac{1}{n}\\
        &= \frac{1}{n^4} + \frac{2^3}{n^4} + \frac{3^3}{n^4} + \cdots +
        \frac{(n - 1)^3}{n^4} + \frac{n^3}{n^4} \\
        & = \frac{1}{n^4}\left(1^3 + 2^3 + 3^3 + \cdots + n^3  \right) \\
        & = \frac{1}{n^4}\left(\dfrac{n^4 + 2n^3 +
            n^2}{4}\right) \text{ by (*)} \\
        & = \boxed{\frac{n^4+2 n^3+n^2}{4 n^4}}
      \end{align*}
      \checkmark\
        When $n = 4$, our formula gives $R_4 = \frac{4^4 + 2 \cdot 4^3 +
          4^2}{4 \cdot 4^4} = \frac{25}{64}$.  When $n = 50$, our formula
        gives $R_{50} = \frac{50^4 + 2 \cdot 50^3 + 50^2}{4 \cdot 50^4} =
        \frac{2601}{10000}$.  These agree with our earlier answers. 

    \end{solution}

  \item
    \begin{problem}[\vfill]
      What is $\displaystyle \lim_{n \to \infty} R_n$?
    \end{problem}

    \begin{solution}
      This is $\displaystyle \lim_{n \to \infty} \frac{n^4+2 n^3+n^2}{4
        n^4}$, which is a $\frac{\infty}{\infty}$ limit.  In the sum $n^4
      + 2n^3 + n^2$, the fastest growing term is $n^4$, so
      \[ \lim_{n \to \infty} \frac{n^4+2 n^3+n^2}{4 n^4} = \lim_{n \to
        \infty} \frac{n^4}{4 n^4} = \boxed{\frac{1}{4}}.\] 
    \end{solution}

  \item
    \begin{problem}[\vfill]
      By definition, $\displaystyle \int_0^1 x^3\,dx$ is $\displaystyle
      \lim_{n \to \infty} R_n$, so you've just computed $\displaystyle
      \int_0^1 x^3\,dx$ exactly.  Does your answer seem reasonable given
      your values for $R_4$ and $R_{50}$ and what you saw in the applet?
    \end{problem}

    \begin{solution}
      Yes, $R_{50} = 0.2601$ is pretty close to $\frac{1}{4}$ (and
      $R_{50}$ should be an overestimate for $\displaystyle \int_0^1
      x^3\,dx$ since $x^3$ is increasing), so our answer seems reasonable.
    \end{solution}

  \end{enumerate}

  \probonly{\emph{As you can see, evaluating definite integrals this way
      is a tremendous amount of work.  In the next few classes, we'll
      develop an alternative (and much less labor-intensive!) way to
      evaluate definite integrals.}}

\end{enumerate}

\psreflection{
Optional reading for next class is \S 22.4, 23.1 in Gottlieb.
}

\end{document}